% Solutions to Chapter 2's exercises, from lectures/L02/appendix/d_solutions.md. No exercise in
% the chapter is a program, so every one is answered, including what the simulation in
% CircuitVerse should show.
%
% Differences in layout only: the derivation in exercise 6 d) is too long for one line of the page,
% so it is set with one step per line, and the shorter route in 6 b) is a display, not inline.

\answerchapter{c2:ch}

\answer{c2:ex:1}{Grindar med tre ingångar}
\needspace{5\baselineskip}
\pt{a}\leavevmode

\begin{narrowtable}{@{}ccc@{\hspace{1.5em}}ccc@{}}
  \thead{A} & \thead{B} & \thead{C} & \thead{AND} & \thead{OR} & \thead{XOR} \\
  \midrule
  0 & 0 & 0 & 0 & 0 & 0 \\
  0 & 0 & 1 & 0 & 1 & 1 \\
  0 & 1 & 0 & 0 & 1 & 1 \\
  0 & 1 & 1 & 0 & 1 & 0 \\
  1 & 0 & 0 & 0 & 1 & 1 \\
  1 & 0 & 1 & 0 & 1 & 0 \\
  1 & 1 & 0 & 0 & 1 & 0 \\
  1 & 1 & 1 & 1 & 1 & 1 \\
\end{narrowtable}

\needspace{5\baselineskip}
\pt{b}\leavevmode
\begin{itemize}
  \item AND är \code{1} när \textbf{alla} tre ingångarna är \code{1}.
  \item OR är \code{1} när \textbf{minst en} ingång är \code{1}.
  \item XOR är \code{1} när ett \textbf{udda} antal ingångar är \code{1}, alltså en eller tre.
\end{itemize}
Den sista raden är den många får fel: XOR med tre ingångar är \code{1} på raden \code{111},
eftersom tre är udda. Den som tänker \enquote{XOR är 1 när ingångarna är olika} missar det, för den
regeln gäller bara med två ingångar.

\pt{c}AND är \code{1} på \textbf{en} rad och OR på \textbf{sju}. AND kräver att alla villkor är
uppfyllda, och det är de bara i en kombination; OR kräver bara ett, och det är uppfyllt i alla
kombinationer utom den där alla är \code{0}. AND och OR är varandras spegelbilder, vilket är precis
vad De Morgan säger.

\answer{c2:ex:2}{Från ord till uttryck}
\pt{a}Med D = \enquote{dörren är öppen}, A = \enquote{larmet är aktiverat} och L = \enquote{larmet
ljuder}: \textbf{\mbox{\code{L = DA}}}.

\pt{b}Med N = \enquote{nivån är låg}, M = \enquote{manuellt stopp är intryckt} och P =
\enquote{pumpen går}: \textbf{\mbox{\code{P = NM'}}}. Ordet \emph{ingen} i beskrivningen är ett NOT.

\pt{c}Med D1, D2 och D3 för de tre dörrarna och H för lampan:
\textbf{\mbox{\code{H = D1 + D2 + D3}}}.

\pt{d}X är \code{1} om maskinen är påslagen och skyddet inte är på plats, eller om nödlarmet är
utlöst. Uttrycket beskriver alltså något som bör leda till ett larm, till exempel ett signalhorn.

\pt{e}\code{AB + C = 0 · 0 + 1 = 1}, men \code{A(B + C) = 0 · (0 + 1) = 0}. De två
uttrycken ger olika svar, så \mbox{\code{AB + C}} får inte läsas som \mbox{\code{A(B + C)}}. AND
binder hårdare än OR, precis som multiplikation binder hårdare än addition.

\answer{c2:ex:3}{Räknelagarna}
\pt{a}\code{A + A = A} (idempotens)
\pt{b}\code{A · A' = 0} (komplement)
\pt{c}\code{A + 1 = 1} (dominans)
\pt{d}\code{(A')' = A} (dubbel invertering)
\pt{e}\code{A + AB = A} (absorption)
\pt{f}\code{A(A + B) = A} (absorption). Eller: \code{A(A + B) = AA + AB = A + AB = A}.
\pt{g}\code{A + A'B = A + B} (förenkling)
\pt{h}\code{AB + A'B = B(A + A') = B · 1 = B} (distributiv, komplement, identitet)

\answer{c2:ex:4}{De Morgan}
\pt{a}\code{(AB)' = A' + B'}
\pt{b}\code{(A + B + C)' = A'B'C'}
\pt{c}\code{(A'B)' = (A')' + B' = A + B'}
\pt{d}\code{(A + B')' = A' · (B')' = A'B}

\ptnote{Vanligt fel}Att i c) och d) bara byta tecknet men glömma att invertera varje variabel, eller
att invertera men glömma att en variabel som redan var inverterad blir oinverterad.

\needspace{5\baselineskip}
\pt{e}\leavevmode

\begin{narrowtable}{@{}cc@{\hspace{1.5em}}ccccc@{}}
  \thead{A} & \thead{B} & \thead{AB} & \thead{(AB)'} & \thead{A'} & \thead{B'} &
    \thead{A' + B'} \\
  \midrule
  0 & 0 & 0 & 1 & 1 & 1 & 1 \\
  0 & 1 & 0 & 1 & 1 & 0 & 1 \\
  1 & 0 & 0 & 1 & 0 & 1 & 1 \\
  1 & 1 & 1 & 0 & 0 & 0 & 0 \\
\end{narrowtable}

\noindent Kolumnerna \code{(AB)'} och \mbox{\code{A' + B'}} är lika på varje rad.

\answer{c2:ex:5}{Från sanningstabell till uttryck}
\pt{a}X är \code{1} på raderna 000, 010, 110 och 111:

\begin{console}
X = A'B'C' + A'BC' + ABC' + ABC
\end{console}

\pt{b}Para ihop de två första termerna (de delar \code{A'C'}) och de två sista (de delar
\code{AB}):

\begin{console}
X = A'C'(B' + B) + AB(C' + C) = A'C' + AB
\end{console}

\pt{c}\code{A'C'} är \code{1} när A = 0 och C = 0: raderna 000 och 010. \code{AB} är \code{1} när
A = 1 och B = 1: raderna 110 och 111. Tillsammans exakt de fyra ettorna i tabellen, och inga andra.

\pt{d}Oförenklat: tre NOT-grindar (för \code{A'}, \code{B'} och \code{C'}), fyra AND-grindar med tre
ingångar och en OR-grind med fyra, alltså \textbf{8 grindar}. Förenklat: två NOT-grindar
(\code{A'} och \code{C'}), två AND-grindar med två ingångar och en OR-grind med två, alltså
\textbf{5 grindar}.

\answer{c2:ex:6}{Algebraisk förenkling}
\pt{a}Använd termen \code{AB} två gånger (idempotens):

\begin{console}
AB + AB' + A'B = (AB + AB') + (AB + A'B) = A(B + B') + B(A + A') = A + B
\end{console}

\needspace{5\baselineskip}
\pt{b}\leavevmode

\begin{console}
(A + B)(A + B') = AA + AB' + AB + BB' = A + A(B' + B) + 0 = A + A = A
\end{console}

\noindent Eller kortare med den andra distributiva lagen baklänges:

\begin{console}
(A + B)(A + B') = A + BB' = A + 0 = A
\end{console}

\pt{c}Samma knep som i a), med \code{SR} två gånger:

\begin{console}
S'R + SR' + SR = (S'R + SR) + (SR' + SR) = R(S' + S) + S(R' + R) = R + S
\end{console}

\needspace{5\baselineskip}
\pt{d}\leavevmode

\begin{textdrawing}
A'B'C + A'BC + ABC = A'C(B' + B) + ABC
                   = A'C + ABC
                   = C(A' + AB)
                   = C(A' + B)
                   = A'C + BC
\end{textdrawing}

\noindent Steget \mbox{\code{A' + AB = A' + B}} är förenklingslagen med A och \code{A'} i omvända
roller. Kontroll: X är \code{1} på raderna 001, 011 och 111, och \mbox{\code{A'C + BC}} är \code{1}
på exakt de raderna.

\answer{c2:ex:7}{Kontroll: analys av ett grindnät}
\needspace{5\baselineskip}
\pt{a}\leavevmode

\begin{console}
P = (A + B)'
Q = BC
X = P + Q = (A + B)' + BC = A'B' + BC
\end{console}

\needspace{5\baselineskip}
\pt{b}\leavevmode

\begin{narrowtable}{@{}ccc@{\hspace{1.5em}}ccc@{}}
  \thead{A} & \thead{B} & \thead{C} & \thead{P = (A + B)'} & \thead{Q = BC} & \thead{X} \\
  \midrule
  0 & 0 & 0 & 1 & 0 & 1 \\
  0 & 0 & 1 & 1 & 0 & 1 \\
  0 & 1 & 0 & 0 & 0 & 0 \\
  0 & 1 & 1 & 0 & 1 & 1 \\
  1 & 0 & 0 & 0 & 0 & 0 \\
  1 & 0 & 1 & 0 & 0 & 0 \\
  1 & 1 & 0 & 0 & 0 & 0 \\
  1 & 1 & 1 & 0 & 1 & 1 \\
\end{narrowtable}

\pt{c}Simuleringen ska ge exakt den här tabellen. De vanligaste orsakerna till en avvikelse är att
NOR-grinden har byggts som en OR (då blir P fel på varje rad), eller att B inte har kopplats till
båda grindarna, eftersom den förgrenar sig i ritningen.

\pt{d}En avvikelse på en enstaka rad beror oftast på ett räknefel i tabellen, eftersom ett fel i
bygget brukar synas på flera rader. Räkna om P och Q för raden, och följ sedan ledningen från varje
ingång i bygget.

\pt{e}X är \code{1} när \textbf{både A och B är \code{0}}, eller när \textbf{både B och C är
\code{1}}.

\answer{c2:ex:8}{Från uttryck till grindnät}
\pt{a}Två NOT-grindar (\code{A'} och \code{C'}), två AND-grindar (\code{A'B} och \code{BC'}) och en
OR-grind: \textbf{5 grindar}.

\needspace{5\baselineskip}
\pt{b}\leavevmode

\begin{console}
X = A'B + BC' = B(A' + C') = B(AC)'
\end{console}

\noindent Det sista steget är De Morgan: \mbox{\code{A' + C' = (AC)'}}. Nätet blir en
\textbf{NAND-grind} med ingångarna A och C, vars utgång går till en \textbf{AND-grind} tillsammans
med B. \textbf{2 grindar} i stället för 5.

\pt{c}En OR-grind för \mbox{\code{A + B}}, en NOT-grind för \code{C'} och en AND-grind: \textbf{3
grindar}.

\pt{d}Båda näten ska vara \code{1} på raderna ABC = 010, 011 och 110, och \code{0} på alla andra.

\answer{c2:ex:9}{Allt av NAND}
\pt{a}Med båda ingångarna kopplade till A blir utgången \mbox{\code{(A · A)' = (A)' = A'}}, eftersom
\mbox{\code{AA = A}}. \textbf{En NAND-grind med ihopkopplade ingångar är en NOT.}

\pt{b}En NAND ger \code{(AB)'}. Inverteras det en gång till blir det \code{AB}. \textbf{Två
NAND-grindar}: en som NAND och en, med ihopkopplade ingångar, som NOT.

\needspace{5\baselineskip}
\pt{c}\leavevmode

\begin{console}
A + B = ((A + B)')' = (A'B')'
\end{console}

\noindent\code{(A'B')'} är en NAND av \code{A'} och \code{B'}, och \code{A'} och \code{B'} är var
sin NAND med ihopkopplade ingångar. Alltså \textbf{tre NAND-grindar}: två som inverterar A och B,
och en tredje som tar emot deras utgångar.

\pt{d}NOT ska ge 1, 0 för A = 0, 1. AND ska bara ge \code{1} för AB = 11. OR ska ge \code{1} för
allt utom AB = 00.

\answer{c2:ex:10}{Från uttryck till kontaktnät}
\pt{a}En strömväg: S1 \textbf{slutande} i serie med S2 \textbf{brytande}, i serie med H1. Lampan
lyser när S1 är nedtryckt och S2 inte är det.

\pt{b}Två parallella grenar mellan +24~V och lampan: den ena med S1 \textbf{slutande}, den andra
med S2 \textbf{slutande} och S3 \textbf{slutande} i serie. Därefter H1.

\pt{c}S1 och S2 \textbf{slutande} parallellt, och den parallellkopplingen i serie med S3
\textbf{brytande} och H1. Samma form som bromsassistentens högra gren i \secref{c2:b6}.

\pt{d}I a)--c) behöver varje knapp \textbf{ett} kontaktelement, eftersom varje variabel förekommer
en gång. I XOR-kopplingen förekommer varje variabel två gånger, en gång rak och en gång inverterad,
så varje knapp behöver \textbf{två}: ett slutande och ett brytande.

\answer{c2:ex:11}{Analys av kontaktnät}
\pt{a}Nät a) är en parallellkoppling av S1 (slutande) och S2 (brytande), i serie med S3:

\begin{console}
H1 = (S1 + S2') · S3
\end{console}

\begin{narrowtable}{@{}cccc@{}}
  \thead{S1} & \thead{S2} & \thead{S3} & \thead{H1} \\
  \midrule
  0 & 0 & 0 & 0 \\
  0 & 0 & 1 & 1 \\
  0 & 1 & 0 & 0 \\
  0 & 1 & 1 & 0 \\
  1 & 0 & 0 & 0 \\
  1 & 0 & 1 & 1 \\
  1 & 1 & 0 & 0 \\
  1 & 1 & 1 & 1 \\
\end{narrowtable}

\pt{b}Nät b) är en seriekoppling av S1 och S2, parallellt med S3 (brytande):

\begin{console}
H2 = S1 · S2 + S3'
\end{console}

\begin{narrowtable}{@{}cccc@{}}
  \thead{S1} & \thead{S2} & \thead{S3} & \thead{H2} \\
  \midrule
  0 & 0 & 0 & 1 \\
  0 & 0 & 1 & 0 \\
  0 & 1 & 0 & 1 \\
  0 & 1 & 1 & 0 \\
  1 & 0 & 0 & 1 \\
  1 & 0 & 1 & 0 \\
  1 & 1 & 0 & 1 \\
  1 & 1 & 1 & 1 \\
\end{narrowtable}

\pt{c}H2 lyser utan S1 och S2 när \textbf{S3 är släppt}: raden 000. S3 är brytande och sitter i en
egen parallell gren, så när S3 inte påverkas finns alltid en sluten väg genom den grenen, vad S1 och
S2 än gör. Faktum är att H2 lyser i alla fyra kombinationer där S3 är släppt.

\answer{c2:ex:12}{Självhållning}
\pt{a}När S1 trycks sluts strömväg 1 genom S2 och S1, och \textbf{spolen K1 drar}. Båda
K1-kontakterna sluts: \textbf{H1 tänds}, och kontakten K1 parallellt med S1 sluts. När S1 släpps går
strömmen genom K1:s egen kontakt i stället, så spolen \textbf{förblir strömsatt} och H1 fortsätter
lysa.

\pt{b}När S2 trycks bryts strömväg 1, \textbf{spolen släpper}, och båda K1-kontakterna öppnar.
\textbf{H1 släcks.} När S2 släpps sluts S2 igen, men S1 och K1 är nu öppna, så spolen förblir
strömlös.

\pt{c}Med S2 nedtryckt är strömväg 1 bruten, vad S1 än gör. \textbf{Stopp vinner.} Kopplingen
kallas därför stoppdominant, och det är så en maskins start- och stoppkrets ska vara.

\pt{d}Med S2 slutande är strömväg 1 bara sluten medan S2 \textbf{hålls nedtryckt}. Att trycka på
start gör ingenting. Trycker man på båda knapparna startar maskinen, och den fortsätter så länge S2
hålls in, men den stannar när man \textbf{släpper} stoppknappen. Stoppknappen fungerar alltså som
en startknapp som måste hållas in, och ingen går att använda för att stoppa maskinen på det sätt en
operatör förväntar sig. Felet är farligt även om det i det här fallet råkar leda till att maskinen
stannar, eftersom ingen kommer att förstå vad som händer.

\pt{e}Med båda knapparna släppta kan H1 vara både tänd och släckt, beroende på om start eller stopp
trycktes senast. Utgången beror alltså inte bara på ingångarna nu, utan på vad som har hänt
tidigare. Kretsen har \textbf{minne}, genom att K1:s kontakt återkopplar spolens läge till dess egen
strömväg, och ett nät med minne är ett sekvensnät, inte ett kombinatoriskt nät.

\answer{c2:ex:13}{Bromsassistenten, felkopplad}
\pt{a}\code{B = (D + S + R)F'} är \code{1} när F = 0 och minst en av D, S och R är \code{1}.
Jämfört med \secref{c2:a10} skiljer sig tabellerna på fyra rader:

\begin{narrowtable}{@{}cccc@{\hspace{1.5em}}cc@{}}
  \thead{D} & \thead{S} & \thead{R} & \thead{F} & \thead{B enligt \secref{c2:a10}} &
    \thead{B felkopplad} \\
  \midrule
  1 & 0 & 0 & 1 & 1 & 0 \\
  1 & 0 & 1 & 1 & 1 & 0 \\
  1 & 1 & 0 & 1 & 1 & 0 \\
  1 & 1 & 1 & 1 & 1 & 0 \\
\end{narrowtable}

\pt{b}På alla fyra raderna är \textbf{D = 1 och F = 1}: föraren trampar på bromspedalen samtidigt
som assistanssystemet rapporterar ett fel. Den felkopplade kretsen \textbf{bromsar inte}. Föraren
ser ett hinder, trampar på bromsen, och ingenting händer, precis i den situation där
assistanssystemet redan har visat att det inte går att lita på.

\pt{c}Sanningstabellen i \secref{c2:a10} skrevs ur kravet, där säkerhetskravet \enquote{föraren kan
alltid bromsa} blev hela den nedre halvan med ettor. En tabell som läses av ur en krets man redan
har byggt visar bara vad kretsen gör, inte vad den borde göra, och skulle ha visat nollor på de fyra
raderna utan att någon reagerade. \textbf{En tabell härledd ur kravet är ett oberoende facit; en
tabell avläst ur kretsen är bara en beskrivning av den.} Det är skälet till att man alltid
förutsäger först och kontrollerar sedan.
